%!TEX encoding = UTF8
%!TEX root = 0-notes.tex

\chapter{Algèbre des polynômes}
\label{chap:polynomes}

\notation{
	Dans tout ce chapitre, $\K$ désigne un corps et $n\in\N$ un entier naturel.
}

Penser d'abord à $\K = \R$ puis à $\C$ : il est possible d'ajouter/soustraitre et de multiplier/diviser les nombres ensembles.
Penser aussi à $\F_p$ ou $\Qsqrt2$.


\section{Anneau $\K[x]$}
\label{sec:Kx}


\dfn{polynôme à coefficients dans $\K$}{
	Un \emphindex{polynôme} à coefficients dans $\K$ est une expression du type
		\[ a_0 + a_1x + a_2x^2 + \dots + a_d x^d, \]
	où $d$ est le \emphindex{degré} du polynôme, les $a_i\in\K$ sont ses \emphindex{coefficients}, et $x$ est une \emphindex{indéterminée} abstraite (pas nécessairement un nombre de $\K$ !).
}{dfn:polynome}

% in depth discussion moved to structures
%\nt{
%	Pour pouvoir évaluer un polynôme, il suffit que l'objet substitué à $x$ appartienne à une $\K$-algèbre
%	(et c'est un bon moyen de se souvenir de la structure de l'algèbre).
%	En effet, 
%	\begin{enumerate}
%		\item il faut pouvoir multiplier $x$ avec lui-même pour calculer ses puissances ; 
%		\item multiplier $x$ avec un élément de $\K$ ; et
%		\item additionner les résultats. 
%	\end{enumerate}
%	Dans ce cas nous verrons $a_0$ comme $a_0\times 1$, où $1$ est l'élément neutre de l'algèbre.
%	
%	Par exemple, il est possible d'évaluer un polynôme en une matrice en voyant $a_0$ comme $a_0 I$.
%}

\notation{
	L'ensemble des polynômes à coefficients dans $\K$ est noté $\K[x]$.
	
	Pour $n\in\N$, l'ensemble des polynômes à coefficients dans $\K$ de degré inférieur ou égal à $n$ est noté $\K_n[x]$.
}

\nomen{
	Le polynôme $a_0 + a_1x + a_2x^2 + \dots + a_d x^d$ est \emphindex{polynôme unitaire} dès que son \emphindex{coefficient dominant} $a_d$ est égal à 1.
}

\notation{
	Le polynôme $a_0 + a_1x + a_2x^2 + \dots + a_d x^d$ est noté $a_i x^i$ en utilisant la \emphindex{convention de sommation d'Einstein} dans laquelle chaque expression dans laquelle un indice $i$ apparaît deux fois sous-entend sa sommation sur toutes les valeurs que prend $i$.
	
	Elle nous permettra de nommer les coefficients sans écrire le polynôme ni définir son degré.
}

\prop{}{
	L'ensemble $\K[x]$ muni des opérations suivantes est un anneau.
	
	Pour $f(x) = f_i x^i$ et $g(x) = g_i x^i$,
	\begin{align*}
		(f+g)(x) = (f_i + g_i)x^i && \et && (f\cdot g)(x) = \left(\sum_{k=0}^{i} f_k g_{i-k}\right)x^i.
	\end{align*}
}{prop:Kx-ring}

\exe{}{
	Montrer que $\K[x]$ est un anneau intègre mais que $\sfrac{\Z}{15\Z}[x]$ n'est pas un anneau intègre.
}{exe:Fpx-integre}{
	
}

\exe{}{
	Montrer que $\K[x]$ est une $\K$-algèbre.
	Dans $\R[x]$, calculer $p(a)$ où $p(x) = x^2 -3x +1$ et $a = 1-x$.
}{exe:Kx-algebra}{

}

\exe{}{
	Donner l'ensemble des éléments inversibles de $\K[x]$.
}{exe:k-star}{
	$\bigpar{ \K[x] }^{\times} = \K^{\times}$.
}

\exe{}{
	Montrer que $\K_n[x]$ est un $\K$-espace vectoriel.
	Donner une base et sa dimension.
	
	Est-ce que $\K_n[x]$ est un anneau ?
}{exe:kn-sur-k}{
	todo
	
	
	Si $n=0$, alors $\K_n[x] = \K$ qui est un anneau.
	Sinon, $\K_n[x]$ n'en est pas un cas la loi $\cdot$ n'est pas interne : $x^n \cdot x^n =x^{2n} \not\in\K_n[x]$.
}

\nt{
	$\K[x]$ est un $\K$-espace vectoriel dans lequel il est possible de multiplier les vecteurs : c'est une $\K$-algèbre.
	Il faut toutefois que la multiplication soit cohérente avec l'addition dans l'espace (distributivité et associativité).
	
	$\K^n$ est un $\K$-espace vectoriel qui n'est pas \emph{a priori} une $\K$-algèbre.
	Elle peut le devenir en munissant l'espace d'un produit interne adéquat.
}

%\subsection{Divisibilité, irréductibilité}

%\dfn{divisiblité dans {$\K[x]$}}{
%	Soient deux polynômes $a, b\in\K[x]$.
%	Le polynôme $a$ divise le polynôme $b$ \underline{dans $\K[x]$} dès que
%		\[ b = a\cdot q, \]
%	où $q\in\K[x]$.
%}{dfn:divisibilite-Kx}
%
%\notation{
%	La relation « $a$ divise $b$ » est notée $a | b$.
%}
%
%\prop{}{
%	Soient $a, b\in\K[x]$. Alors
%		\begin{align*}
%			a | b && \iff && (b) \subseteq (a).
%		\end{align*}
%}{prop:incl-is-div}

\exe{}{
	Soient $f, g\in\K[x]$ deux polynômes.
	Montrer que si $f|g$ et $g|f$, alors $f = c \cdot g$ avec $c\in\K$.
	
	En déduire que si $f$ et $g$ sont unitaires, alors $f=g$.
}{exe:f-div-g-div-f}{
	todo
}

\nt{
	Il est important de préciser dans quel anneau polynomial la division s'opère.
	Il peut s'avérer utile d'agrandir le corps dans lequel nous travaillons pour étudier la divisibilité.
	Par exemple, tout polynôme de $\R[x]$ admet une racine sur $\C$, et donc voir ce même polynôme comme appartenant à $\C[x]$ permet de dégager un diviseur.
	Le polynôme réel est vu comme un polynôme complexe par \emphindex{plongement} de $\R$ dans $\C$.
}

\exe{}{
	Soient $a, b\in\R[x]$ deux polynômes à coefficients réels.
	Montrer que si $a|b$ dans $\R[x]$, alors $a|b$ dans $\C[x]$.
}{exe:divis-R-to-C}{
	todo
}

\exe{, difficulty=1}{
	Soient $a, b\in\R[x]$ deux polynômes à coefficients réels.
	Montrer que si $a|b$ dans $\C[x]$, alors $a|b$ dans $\R[x]$.
}{exe:divis-C-to-R}{
	Notons $b = aq +r, \deg(r)<\deg(a)$, la division euclidienne de $b$ par $a$ dans $\R[x]$.
	Or, par hypothèse, nous avons $b = az$, avec $z\in\C[x]$.
	Dans $\C[x]$ toujours, nous obtenons donc $az = aq + r \iff r = a(z-q)$.
	Or comme $\deg(r)<\deg(a)$, il suit nécessairement que $z-q=r=0$.
	Par conséquent, $b$ divise $a$ dans $\R[x]$.
}

%\dfn{polynôme irréductible}{
%	Un polynôme $f\in\K[x]$ est un \emphindex{polynôme irréductible} dès que
%		\begin{enumerate}
%			\item $f$ n'est pas inversible ; et
%			\item si $f$ se scinde $f=ab$ dans $\K[x]$, alors soit $a$ soit $b$ est inversible.
%		\end{enumerate}
%}{dfn:irreductible-Kx}

%\nt{
%	La définition parle de « polynôme inversible » au lieu de « constante » de sorte qu'elle se généralise aux anneaux $R[x]$ où $R$ est un anneau, non nécessairement un corps.
%	Dans $\K[x]$, les éléments inversibles sont exactement les constantes (hormis 0), comme vu à l'exercice \ref{exe:k-star}.
%}

\exe{}{
	Le polynôme $f(x) = x^4+x^3+x^2+x+1$ est-il irréductible dans $\F_5[x]$ ?
}{exe:irreductible}{

}

\exe{}{
	Soit $f\in\Z[x]$ unitaire et $p$ un premier. 
	Montrer que si $(f\mod p)$ est irréductible sur $\F_p[x]$, alors $f$ est irréductible dans $\Z[x]$.
	
	Montrer que le résultat devient faux si $f$ n'est pas supposé unitaire.
	
	Est-ce que la réciproque est vraie ?
}{exe:irreductible-Z-to-Fp}{
	Par contraposition, supposons $f$ réductible dans $\Z[x]$ : $f = ab$, avec $a, b\in\Z[x]$ non constants.
	Comme $f$ est unitaire, $a$ et $b$ le sont aussi.
	En prenant l'égalité modulo $p$, nous obtenons donc
		\[ (f \mod p) = (a \mod p)(b \mod p), \]
	égalité dans $\F_p[x]$.
	Les facteurs $(a \mod p)$ et $(b\mod p)$ sont unitaires dans $\F_p[x]$ est donc non constants.
	
	Si $f$ n'est pas supposé unitaire, alors le polynôme de Clamme $\text{Cl}(x) = px^2 + x$ fournit un contre-exemple.
	
	La réciproque est fausse, car $x^2 +1$ est irréductible dans $\Z[x]$ mais se réduit dans $\F_2[x]$ comme $x^2 + 1 = (x+1)^2$.
}

\exe{difficulty=1}{
	Soit $p\in\N$ un premier impair.
	\begin{enumerate}
		\item
		Montrer que le polynôme $ax^2 + bx + c \in \Z[x]$ avec $p\nmid a$ est réductible dans $\F_p[x]$ si et seulement si le discriminant $\Delta = b^2 - 4ac$ est un carré parfait modulo $p$.
		\item 
		En déduire un polynôme de $\Z[x]$ irréductible dans $\F_5[x]$ mais réductible dans $\F_7[x]$.
		\item
		Montrer que $x^2 + x + 1$ est irréductible dans $\F_2[x]$.
	\end{enumerate}
}{exe:discr-Fp}{
	Le polynôme est de degré 2 et est donc réductible si et seulement s'il admet une racine dans $\F_p$.
	Les équations suivantes sont équivalentes et admettent une solution si et seulement si $\Delta$ est un carré parfait.
	\begin{align*}
		ax^2 + bx + c = 0 \\
		4a^2x^2 + 4abx + 4ac = 0 \\
		(2ax+b)^2 - b^2 + 4ac = 0 \\
		(2ax+b)^2 = \Delta
	\end{align*}
	Nous avons utilisé ici que $p\neq2$ pour avoir que $4$ est inversible dans $\F_p$, et que $p\nmid a$ pour que $a$ soit inversible dans $\F_p$.

	Les carrés parfaits de $\F_5$ sont $\bigset{ 0 ; 1 ; 4 }$.
	Ceux de $\F_7$ sont $\bigset{ 0 ; 1 ; 4 ; 2 }$.
	Obtenir $\Delta = 2$ est donc notre unique chance et le polynôme $x^2 - 2$ un candidat évident.
	Celui-ci est irréductible dans $\F_5[x]$ mais se décompose comme $(x-3)(x+3)$ dans $\F_7[x]$.
	
	Il existe des tables de polynômes irréductibles pour chaque corps fini, dans lesquelles nous pouvons trouver 
	\begin{itemize}[label=$\bullet)$]
		\item $x^2 + x + 1$, de discriminant 2 dans $\F_5$ (irréductible) et 3 dans $\F_7$ (irréductible) ;
		\item $x^2 + x + 2$, de discriminant 3 dans $\F_5$ (irréductible) et 0 dans $\F_7$ (réductible) ;
		\item $x^2 + 4x + 1$, de discriminant 2 dans $\F_5$ (irréductible) et 5 dans $\F_7$ (irréductible) ;
		\item etc...
	\end{itemize}
}


%\nt{
%	Le \emphindex{lemme de Gauss (polynômes)} énonce le fait suivant :
%	si $a\in\Q[x]$ divise $b\in\Z[x]$, deux polynômes unitaires, dans $\Q[x]$, alors $a\in\Z[x]$ est à coefficients entiers et $a$ divise $b$ dans $\Z[x]$.
%	
%	Ce résultat semble facile à démontrer en utilisant la même stratégie de preuve qu'à l'exercice \ref{exe:divis-C-to-R}, mais la division euclidienne sort de $\Z[x]$ en général : le polynôme $3(x+1)$ divise $x^2 - 1$ dans $\Q[x]$ mais pas dans $\Z[x]$. (Notons que le premier n'est pas unitaire !)
%}

%\dfn{PGCD}{
%	Soient $a, b\in\K[x]$ deux polynômes à coefficients dans $\K$.
%	Un plus grand diviseur commun (PGCD) et $a$ et $b$ dans $\K[x]$ est un polynôme $d\in\K[x]$ de degré maximal vérifiant
%	\begin{enumerate}
%		\item $d$ divise $a$ dans $\K[x]$ ; et
%		\item $d$ divise $b$ dans $\K[x]$ ; et
%		\item $d$ est unitaire.
%		\item $(\star)$ les diviseurs de $a$ et de $b$ 
%	\end{enumerate}
%}{dfn:pgcd}


\subsection{L'anneau $\K[x]$ est euclidien}

\dfn{degré}{
	Soit $f(x) = a_0 + a_1x + \dots + a_dx^d\in\K[x]$ un polynôme à coefficients dans $\K$.
	Le \emphindex{degré} de $f$ est défini par
	\begin{itemize}
		\item
		$\deg(f) = d$ si $a_d \neq 0$ ; ou
		\item
		$\deg(f) = \minfty$ si $f = 0$ est le polynôme nul.
	\end{itemize}
}{dfn:degre}

\nt{
	La convention $\deg(0) = \minfty$ est utile pour préserver le premier point de la proposition \ref{prop:deg-div} suivante dans le cas $g=0$.
}

\prop{}{
	Soient $f, g\in\K[x]$. Alors
	\begin{enumerate}
		\item
		$\deg(fg) = \deg(f) + \deg(g)$
		%\item
		%$\deg(f+g) \leq \max\{\deg(f) ; \deg(g)\}$
		\item
		Si $f|g$ et $g\neq0$, alors $\deg(f) \leq \deg(g)$
	\end{enumerate}
}{prop:deg-div}

\pf{}{
	Si $f = 0$ ou $g=0$, alors il n'y a rien à démontrer.
	
	Sinon, notons $f(x) = f_i x^i$ et $g(x) = g_i x^i$ de degrés $m$ et $n$ respectivement.
	D'après le corollaire \ref{cor:corps-integre}, le corps $\K$ est intègre.
	Ainsi, $(fg)(x)= f_mg_nx^{m+n} + \dots$ et $\deg(fg) = m+n$ car $f_m g_n \neq 0$.
}

\exe{}{
	Quand est-ce que l'identité $\deg(f+g) = \max\{\deg(f) ; \deg(g)\}$ est vraie ?
}{exe:deg-sum}{
	Il faut et il suffit que la somme des coefficients de degré $\max\{\deg(f) ; \deg(g)\}$ soit non nulle.
	Si $g=-f$ par exemple, même l'inégalité $\deg(f+g) \leq \max\{\deg(f) ; \deg(g)\}$ est fausse.
}

\thm{division euclidienne}{
	Soient $a, b\in\K[x]$ deux polynômes à coefficients dans $\K$ avec $a\neq0$.
	Il existe deux uniques polynômes $q, r\in\K[x]$ vérifiant
	\begin{enumerate}
		\item $b = aq + r$ ; et
		\item $\deg(r) < \deg(a)$
	\end{enumerate}
	C'est la \emphindex{division euclidienne} de $b$ par $a$.
}{thm:div-euclidienne}

\nt{
	Il suit naturellement que $a$ divise $b$ dans $\K[x]$ si et seulement si le reste $r$ de la division euclidienne de $b$ par $a$ est nul.
	
	Remarquons que si $a=0$, alors la division euclidienne de $b$ par $a$ soit n'existe pas ($b\neq0$), soit n'est pas unique ($b=0$).
}

\exe{}{
	Démontrer le théorème \ref{thm:div-euclidienne}, sans oublier l'unicité de $q$ et $r$.
	Est-ce qu'il est toujours vrai si on remplace le corps $\K$ par un anneau $R$ ?
}{exe:thm:div-euclidienne}{
	\underline{Existence.}

	Si $\deg(b) < \deg(a)$, alors il suffit de prendre $q=0$ et $r=b$.
	Sinon, $\deg(b) \geq \deg(a)$ et nous construisons $q$ en annulant le degré le plus haut de $b$ itérativement.
	
	\begin{enumerate}
		\item
		Initialiser $q = 0$.
		\item
		Tant que $\deg(b-aq) \geq \deg(a)$, 
		\item
		Ajouter un multiple de $x^{\deg(b-aq) - \deg(a)}$ à $q$ de telle sorte que le degré de $b-aq$ diminue d'au moins 1.
		\item
		Fin Tant que
		\item
		Renvoyer $q$ et $r=b-aq$.
	\end{enumerate}

	\underline{Unicité.}
	
	Si $b = aq + r = aq' + r'$ avec $\deg(r), \deg(r') < \deg(a)$, alors
		\[ a(q-q') = r' - r, \]
	et donc $q-q'=0$ car $\deg(r'-r) < \deg(a)$ et que $a\neq0$.
	Il suit que $q=q'$ et donc que $r=r'$.
}

\nt{
	Le concept d'\emphindex{anneau euclidien} se généralise sans peine en remplaçant le degré par une norme sur l'anneau $R$ qui prend ses valeurs dans $\N$.
}

\exe{difficulty=1}{
	Soient $a, b \in\Z[i]$.
	Montrer qu'il existe $q, r\in\Z[i]$ vérifiant $b = aq+r$ et $|r| < |a|$.
}{exe:Zi-euclidien}{
	todo
}

\nt{
	L'interprétation de l'exercice \ref{exe:Zi-euclidien} est la suivante.
	L'anneau des \emphindex{entiers de Gauss} $\Z[i]$ est euclidien avec la norme $|\cdot|$ car c'est un réseau de $\C$ suffisamment fin pour que tout élément de $\Q[i]$ admette un élément de $\Z[i]$ dans sa boule unité.
	
	Ainsi, en élargissant le réseau, il est naturel que l'anneau perde sa division euclidienne.
	Il s'avère alors que $\Z\left[i\sqrt2\right]$ est toujours euclidien, mais pas $\Z\left[i\sqrt3\right]$.
}

\exe{difficulty=2}{
	Montrer que la division euclidienne de $1+i\sqrt3$ par 2 n'existe pas dans $\Z[i\sqrt3]$.
}{exe:Zsqrt3i-noneuclidien}{
	todo
}

\exe{}{
	Soit $f\in\K[x]$ un polynôme non constant.
	Montrer que si $r\in\K$ est une racine de $f$, c'est-à-dire si $f(r) = 0$, alors $x-r$ divise $f$ dans $\K[x]$.
}{exe:root-factor}{
	todo
}

\exe{difficulty=2}{
	Soit $f\in\K[x]$ un polynôme non constant de degré $n\in\N$.
	Montrer l'équivalence suivante.
	\begin{center}
	\begin{tabular}{ccc}
		$\Bigpar{ \text{ $f$ est réductible } }$ & $\iff$ & $\Bigpar{\text{ il existe $a, b\in\K[x]$ non nuls avec $\deg(a), \deg(b) < n$ tels que $f|(ab)$ }}$
	\end{tabular}
	\end{center}
}{exe:red-div-prod}{
	\underline{Direction $\implies$.}
	
	Si $f = ab$ avec $a, b$ non constants, alors $f$ divise trivialement $ab$ et $\deg(a) + \deg(b) = n$ implique que $\deg(a), \deg(b) < n$.
	
	\underline{Direction $\impliedby$.}
	
	Considérons l'ensemble $E$ des couples $(a, b)$ de polynômes \underline{non nuls} vérifiant $\deg(a), \deg(b) < n$ et $f|(ab)$.
	Cet ensemble n'est pas vide par hypothèse.
	Choisissons un couple de $E$ pour lequel $\deg(a)$ est minimal et montrons que $a$ divise $f$.
	
	La division euclidienne de $f$ par $a$ dans $\K[x]$ s'écrit
		\[ f = aq + r, \]
	avec $\deg(r) < \deg(a)$.
	
	Comme $r = f -aq$, le polynôme $f$ divise $rb = rf - (ab)q$.
	Si $r\neq0$, l'appartenance $(r,b) \in E$ contredit le choix de $a$.
	Il suit donc que $r = 0$ et nous avons trouvé un diviseur propre de $f$.
	En effet, $a$ ne peut pas être constant car $\deg(b) < n$ et $\deg(ab) \geq n$.
}

\subsection{L'anneau $\K[x]$ est principal}


\dfn{anneau principal}{
	Un anneau intègre $R$ est un \emphindex{anneau principal} dès que tout idéal de $R$ est généré par un unique élément.
}{dfn:anneau-principal}


\exe{difficulty=1}{
	Montrer que $\K[x]$ est un anneau principal : tout idéal de l'anneau est généré par l'élément de degré minimal.
}{exe:Kx-is-PID}{
	todo
}

\exe{}{
	Montrer que l'anneau $\Z[x]$ n'est pas principal : il n'existe pas de polynôme $f\in\Z[x]$ tel que $(2, x) = (f)$.
}{exe:Zx-non-principal}{
	todo
}

\dfn{PGCD}{
	Soient $a, b\in \K[x]$ deux polynômes à coefficients dans $\K$.
	Un plus grand diviseur commun (PGCD) et $a$ et $b$ dans $\K[x]$ est un élément $d\in\K[x]$ vérifiant
	\begin{enumerate}
		\item $d$ divise $a$ et $b$ ; et
		\item les diviseurs de $a$ et de $b$ divisent $d$ ; et
		\item $d$ est unitaire.
	\end{enumerate}
}{dfn:pgcd}

\nt{
	L'existence du PGCD n'est \apriori~pas claire : il se pourrait que, pour chaque candidat de PGCD $d$, il existe un diviseur de $a$ et de $b$ ne divisant pas $d$.
	Nous verrons dans la suite que le PGCD existe bien dans $\K[x]$.
	%L'anneau $\Z[\sqrt5i]$ est un exemple classique d'anneau sans PGCD.
	
	L'unicité, elle, est garantie par l'exercice \ref{exe:unicite-pgcd} suivant.
}

\exe{}{
	Montrer que le PGCD, s'il existe, est unique.
}{exe:unicite-pgcd}{
	Soient $d, d'$ deux PGCD de $a$ et $b$ dans $\K[x]$.
	Comme $d$ divise à la fois $a$ et $b$, il divise $d'$.
	Symétriquement, $d'$ divise $d$.
	Par conséquent, $d$ et $d'$ sont multiples l'un de l'autre, d'après l'exercice \ref{exe:f-div-g-div-f}.
	Leur caractère unitaire permet de conclure que $d=d'$.
}


\dfn{coprimalité}{
	Deux polynômes $f, g\in\K[x]$ sont \emphindex{premiers entre eux} dès que $\pgcd(f, g) = 1$.
}{dfn:coprimalite}

\exe{}{
	Montrer que si $d|a$ et $d|b$ dans $\K[x]$, alors $d$ divise toutes les combinaisons linéaires $au + bv$ avec $u, v\in\K[x]$ également.
}{exe:d-div-comb-lin}{
	todo
}

\exe{}{
	Soient $f\in\K[x]$ irréductible et $g\in\K[x]$ non nul tel que $\deg(g) < \deg(f)$.
	Montrer que $\pgcd(f,g) = 1$.
}{exe:irr-coprime}{
	todo
}

\exe{}{
	Soit $f\in\K[x]$ irréductible et $a\in\K[x]$.
	Montrer que $\pgcd(f,a) = 1$ ou $f$.
}{exe:gcd-irr}{
	todo
}


\thm{de Bézout}{
	Soient $a, b\in\K[x]$ deux polynômes à coefficients dans $\K$.
	% et $d\in\K[x]$ un PGCD de $a$ et $b$.
	
	Alors le PGCD de $a$ et $b$ existe est égal à
		\[ \pgcd(a, b) = \arg\min\bigset{ \deg(p) \tq p = au + bv \text{ avec } u, v \in \K[x], \et p \text{ unitaire non nul} }. \]
}{thm:bezout-K}

\pf{}{
	Notons qu'un argument du minimum n'est pas nécessairement unique \apriori.
	Montrons donc d'abord qu'il l'est en prenant $p$ et $p'$ deux arguments du minimum.
	Les deux polynômes $p$ et $p'$ sont donc combinaisons linéaires de $a$ et $b$, non nuls, unitaires, et de degré minimal égal à $d\in\N$.
	La différence $p-p'$ est de degré strictement inférieur à $p$ et est combinaison linéaire de $a$ et $b$.
	Par minimalité de $p$, cette différence est multiple de 0 et est donc nulle : $p=p'$.
	
	Notons ainsi $p$ l'unique argument du minimum et montrons qu'il est un PGCD de $a$ et $b$,
	c'est-à-dire qu'il vérifie les trois conditions suivantes.
	\begin{enumerate}
		\item $p$ divise $a$ et $b$ dans $\K[x]$ ; et
		\item les diviseurs de $a$ et de $b$ divisent $p$ ; et
		\item $p$ est unitaire ; 
	\end{enumerate}
	\begin{enumerate}[label=\underline{Condition \arabic*.}]
		\item
		%La condition 1 se démontre classiquement comme suit.
		En prenant $a = pq + r, \deg(r) < \deg(p)$ la division euclidienne de $a$ par $p$ dans $\K[x]$, nous obtenons que $r = a-pq$ est combinaison linéaire de $a$ et $b$ et est donc nul par minimalité du degré de $p$.
		Il suit que $p$ divise $a$ dans $\K[x]$, et que $p$ divise également $b$ de façon analogue.
		
		\item
		D'après l'exercice \ref{exe:d-div-comb-lin}, un diviseur $d\in\K[x]$ de $a$ et de $b$ divise $p$ car c'est une combinaison linéaire de $a$ et de $b$.
		%La condition 2 est donc vérifiée.
		
		\item		
		%La condition 3 est vérifiée par construction de $p$.
		Vérifiée par construction de $p$.
	\end{enumerate}
}

\exe{}{
	Soient $a, b\in\K[x]$ deux polynômes unitaires.
	Montrer que si $a$ divise $b$ dans $\K[x]$ et que $\deg(b) \leq \deg(a)$, alors $a=b$.
}{exe:div-deg-eq}{
	todo
}

\exe{}{
	Dans cet exercice, nous considérons une autre définition du PGCD, dont l'existence est claire, mais l'unicité ne l'est pas.
	
	Soient $a, b\in\K[x]$ deux polynômes à coefficients dans $\K$.
	Un plus grand diviseur commun prime (PGCD') et $a$ et $b$ dans $\K[x]$ est un polynôme $d\in\K[x]$ de degré maximal vérifiant
	\begin{enumerate}
		\item $d$ divise $a$ et $b$ dans $\K[x]$ ; et
		\item $d$ est unitaire.
	\end{enumerate}
	
	Montrer que $\pgcd'(a,b)$ est unique et qu'il est égal à $\pgcd(a,b)$.
}{exe:pgcd-other-def}{
	Notons $p = \pgcd(a,b)$ et $p'=\pgcd'(a,b)$.
	
	D'abord, comme $p'$ divise $a$ et $b$ dans $\K[x]$, alors $p'$ divise $p$, par définitions.
	
	Ensuite, comme $p$ divise $a$ et $b$ dans $\K[x]$, alors $\deg(p) \leq \deg(p')$, à nouveau par définitions.
	
	Il suit donc que $p=p'$ en utilisant leur caractère unitaire et l'exercice \ref{exe:div-deg-eq}.
}

\cor{}{
	Soient $a, b \in\K[x]$. Alors
		\[ (a, b) = \bigpar{ \pgcd(a, b) }. \]
	%En particulier, si $d\in\K[x]$ divise à la fois $a$ et $b$, alors $d$ divise $\pgcd(a, b)$.
	En particulier, il existe $u, v\in\K[x]$ tels que
		\[ \pgcd(a,b) = au + bv. \]
}{cor:pgcd-ideal}

\exe{}{
	Montrer le corollaire \ref{cor:pgcd-ideal}.
}{exe:cor:pgcd-ideal}{
	todo
}

\exe{difficulty=1}{
	Soient $\K$ un corps de caractéristique nulle, $f\in\K[x]$ quelconque, et $g\in\K[x]$ irréductible.
	Supposons que $g$ divise $f$ et $f'$, la dérivée formelle de $f$.
	Montrer que $g^2 | f$.
	
	Donner un contre-exemple en caractéristique finie.
}{exe:g2-div-f}{
	Comme $g$ divise $f$, nous avons $f=qg$ et $f' = q'g + qg'$.
	De plus, $g$ divise $f'$ et donc $f'-q'g = qg'$.
	Or, $g$ est irréductible, donc il est premier avec $g'$ car de degré inférieur et non nul (on utilise ici que $\K$ est de caractéristique nulle).
	Le lemme d'Euclide \ref{lem:Euclide} 
	
	En caractéristique $p$, le polynôme $x^p$ divise lui-même et sa dérivée, mais son carré ne le divise pas.
}

%\nt{
%	L'anneau $\K[x]$ est donc un \emphindex{anneau principal} : l'élément de degré (ou norme, plus généralement) minimal de chaque idéal génère entièrement cet idéal.
%}

%\nt{
%	Une définition plus abstraite mais plus générale du PGCD est donc la suivante, et s'applique à tous les anneaux principaux.
%	Pour $a, b\in\R$ deux éléments d'un anneau $R$ principal, nous avons
%		\[ (a, b) = \bigpar{ \pgcd(a, b) }, \]
%	ce qui définit $\pgcd(a, b)$ à élément inversible de $R$ près.
%	Dans notre cas, les éléments inversibles de $\K[x]$ sont les constantes non nulles, et le caractère unitaire du PGCD permet de le rendre unique.
%}

\subsection{L'anneau $\K[x]$ est factoriel}

% eh, not sure it's needed
%\lem{de Gauss (arithmétique)}{
%	Soit $a, b, c\in\K[x]$ trois polynômes à coefficients dans $K$. Les divisions sont comprises dans $\K[x]$.
%	
%	Si $a$ divise le produit $bc$, et que $a$ et $b$ sont premiers entre eux, alors $a$ divise $c$.
%}{lem:Gauss-arithm}

\lem{d'Euclide}{
	Soit $p\in\K[x]$ un polynôme irréductible, et $a, b\in\K[x]$ deux polynômes. 
	
	Si $p|(ab)$, alors $p|a$ ou $p|b$.
}{lem:Euclide}

\pf{}{
	Si $p$ divise $a$, il n'y a rien à démontrer.
	Supposons donc que $p$ ne divise pas $a$ et montrons que $p$ divise $b$.
	
	Le PGCD $d = \pgcd(p, a)$ divise $p$, qui est irréductible. Ainsi, $d = 1$ ou $d=p$, la deuxième alternative n'étant pas possible car $p$ ne divise pas $a$ par hypothèse.
	
	Alors $\pgcd(p,a) = 1$ et le théorème de Bézout \ref{thm:bezout-K} implique l'existence de $u, v\in\K[x]$ vérifiant
		\[ pu + av = 1. \]
	En multipliant par $b$, nous obtenons
		\[ p(ub) + (ab)v = b, \]
	divisible par $p$ car $p|(ab)$, ce qui conclut.
}

\exe{}{
	Montrer le \emphindex{lemme de Gauss (arithmétique)} suivant.
	
	Soit $a, b, c\in\K[x]$ trois polynômes.
	Si $a|(bc)$ et $\pgcd(a,b) = 1$, alors $a|c$.
}{exe:lem:Gauss}{
	Or, come $\pgcd(a,b)=1$, il existe $u, v\in\K[x]$ tels que
		\[ au + bv = 1. \]
	En multipliant par $c$, nous obtenons que
		\[ c = a(uc) + (bc)v \]
	est divisible par $a$.
}
 
\cor{de l'exercice \ref{exe:lem:Gauss}}{
	Soit $f\in\K[x]$ un polynôme quelconque. Les divisions et PGCD sont compris dans $\K[x]$.
	
	Si $p, q\in\K[x]$ divisent tous les deux $f$ et que $\pgcd(p,q) = 1$, alors leur produit $pq$ divise également $f$.
}{cor:Gauss-arithm}

\nt{
	Ce corollaire est utile pour justifier immédiatement que si $p$ admet plusieurs racines distinctes $r_1, r_2, \dots, r_k$, alors le produit $\prod_{i=1}^k (x-r_i)$ divise $p$ dès qu'il appartient à l'anneau.
}

\exe{}{
	Montrer le corollaire \ref{cor:Gauss-arithm}.
}{exe:cor:Gauss-arithm}{
	Comme $p$ divise $f$, nous avons $f=ap$ avec $a\in\K[x]$.
	Par suite, $q$ divise le produit $ap$ et est premier avec $p$ par hypothèse.
	Le lemme de Gauss (arithmétique) de l'exercice \ref{exe:lem:Gauss} conclut que $q$ divise $a$.
	Dès lors, $a = bq$ et donc $f = bqp$ comme voulu.
}

\exe{difficulty=1}{
	Soient $a, b\in\K[x]$.
	Le PPCM est défini comme le polynôme unitaire vérifiant
		\[ (a) \cap (b) = \bigl( \ppcm(a, b) \bigr). \]
	Son existence est garantie car $(a) \cap (b)$ est un idéal de $\K[x]$ qui est principal.
	\begin{enumerate}
		\item
		Montrer que le PPCM est l'unique polynôme unitaire $p\in\K[x]$ vérifiant
		\begin{enumerate}
			\item $a|p$ et $b|p$ ; et
			\item si $a|d$ et $b|d$, alors $p|d$.
		\end{enumerate}
		\item
		Montrer que $ab = \ppcm(a,b)\cdot\pgcd(a,b)$.
	\end{enumerate}
}{exe:ppcm}{
	$au + bv = \pgcd(a,b)$ \\
	$aq = d$ \\
	$du + bqv = \pgcd(a,b)q$ \\
	$b|\pgcd(a,b)q$\\
	$db = p \pgcd(a,b)q$\\
	$d = p q'$\\
	$p|d$.
}

\nt{
	Les polynômes irréductibles sont les équivalents des nombres premiers.
	Il est donc naturel de vouloir trouver un analogue au théorème fondamental de l'arithmétique dans l'anneau des polynômes.
	Pour cela, il convient de se demander : quels polynômes sont irréductibles ? 
	La réponse à cette question dépend du corps choisi. 
	
	Sur $\C$, seuls les polynômes de degré 1 sont irréductibles. 
	Sur $\R$, des polynômes de degré 2 peuvent l'être, comme $x^2 + 1$, mais pas de degré 3 car le théorème des valeurs intermédiaires implique l'existence d'une racine réelle.
	Sur $\Z$ et les corps finis, il existe des polynômes irréductibles de degrés aussi grands que voulu.
}

\thm{}{
	Soit $f\in\K[x]$ un polynôme quelconque.
	Alors $f$ se décompose dans $\K[x]$ de façon unique comme produit de facteurs irréductibles.
	L'unicité est à permutation des facteurs et à multiplication de chaque facteur par une unité près.
}{thm:decomp-Kx}

\nt{
	Il est possible de décomposer $x^2 - 1$ dans $\Q[x]$ comme
		\[ x^2 -1 = \left( \frac12 x - \frac12 \right) (2x + 2) = (x+1)(x-1), \]
	qui sont deux décompositions identiques à permutation des facteurs et à multiplication de chaque facteur par une constante près.
}

\pf{}{
	L'existence est claire : un polynôme est soit irréductible, soit il se décompose en deux facteurs de degrés strictement inférieurs. 
	Ce processus se termine nécessairement, les degrés étant des entiers positifs ou nuls.
	
	Considérons désormais deux décompositions en produits de facteurs irréductibles de $f$.
	Il suffit de montrer que chaque facteur de la première décomposition apparaît dans la seconde.
	Par symétrie, chaque facteur de la seconde apparaît dans la première, et ceci conclut.
	
	Soit donc $p$ un facteur de la première décomposition. 
	Comme $p$ divise la deuxième décomposition, le lemme d'Euclide \ref{lem:Euclide} implique que $p$ divise un de ses facteurs irréductibles.
	$p$ est donc égal à ce facteur.
}

\exe{}{
	Montrer qu'un polynôme $f$ de $\K[x]$ de degré 2 ou 3 est réductible si et seulement s'il admet une racine dans $\K$.
}{exe:irr-deg23}{
	\underline{Direction $\implies$.}
	
	Si se scinde comme $f=ab$ avec ni $a$ ni $b$ constants, alors soit $a$ soit $b$ est de degré 1 car $\deg(a) + \deg(b) = \deg(f)$.
	Le facteur de degré 1 implique donc l'existence d'une racine dans le corps.
	
	 
	\underline{Direction $\impliedby$.}
	
	Si $r\in\K$ vérifie $f(r) = 0$, alors $x-r$ divise $f$ dans $\K[x]$ d'après l'exercice \ref{exe:root-factor}.
	Le polynôme $f$ est donc réductible.
}

\exe{}{
	Trouver la décomposition de chaque polynôme suivant dans les anneaux $\C[x] ; \R[x] ; \Q[x] ; \F_3[x]$.
	\begin{multicols}{4}
	\begin{enumerate}
		\item $x^2 + 2$
		\item $x^2 - 1$
		\item $x^2 - 2$
		\item $x^3 + x^2 + 2x + 2$
	\end{enumerate}
	\end{multicols}
}{exe:UFD-examples}{
	Chaque polynôme est réductible si et seulement si nous pouvons trouver une racine dans le corps, d'après l'exercice \ref{exe:irr-deg23}.
	\begin{enumerate}
		\item $x^2 + 2$.
		\begin{enumerate}[label=Dans~, labelsep=0pt,]
			\item$\C[x]$ :  $x^2 + 2 = (x + \sqrt2i)(x-\sqrt2i)$.
			\item$\R[x]$ : $x^2 + 2$ est irréductible.
			\item$\Q[x]$ : $x^2 + 2$ est irréductible.
			\item$\F_3[x]$ : $x^2 + 2 = x^2 - 1 = (x+1)(x-1)$.
		\end{enumerate}
		\item $x^2 - 1 = (x+1)(x-1)$ dans chaque anneau.
		\item $x^2 - 2$.
		\begin{enumerate}[label=Dans~, labelsep=0pt,]
			\item$\C[x]$ :  $x^2 - 2 = (x+\sqrt2)(x-\sqrt2)$.
			\item$\R[x]$ : $x^2 - 2 = (x+\sqrt2)(x-\sqrt2)$.
			\item$\Q[x]$ : $x^2-2$ est irréductible.
			\item$\F_3[x]$ : $x^2-2$ est irréductible.
		\end{enumerate}
		\item $x^3 + x^2 + 2x + 2$ dont $-1$ est une racine et qui se factorise donc par $x+1$ dans chaque anneau.
		\begin{enumerate}[label=Dans~, labelsep=0pt,]
			\item$\C[x]$ :  $x^3 + x^2 + 2x + 2 = (x+1)(x^2+2) = (x+1)(x+\sqrt2i)(x-\sqrt2i)$
			\item$\R[x]$ : $x^3 + x^2 +2x +2 = (x+1)(x^2 + 2)$
			\item$\Q[x]$ : $x^3 + x^2 + 2x + 2 = (x+1)(x^2 + 2)$
			\item$\F_3[x]$ : $x^3 + x^2 + 2x + 2 = (x+1)(x^2 -1) = (x+1)^2(x-1)$
		\end{enumerate}
	\end{enumerate}
}


\exe{}{
	Montrer que l'anneau $\sfrac{\Z}{8\Z}[x]$ n'est pas factoriel en décomposant le polynôme $x^2 -1$ de deux façons différentes.
}{exe:Zo8Zx-non-factoriel}{
	L'élément 1 admet quatre racines carrés dans $\Z/8\Z$ !
	Celles-ci sont $\pm1$ et $\pm3$.
	Ainsi $x^2 - 1 = (x+1)(x-1) = (x+3)(x-3)$ dans l'anneau.
}

\section{L'anneau $\Z[x]$}
\label{sec:Zx}

\nt{
	L'anneau $\Z[x]$ n'a pas de division euclidienne interne car $\Z$ n'est pas un corps.
	En effet, comment procéder à la division de $x^5 - 3x + 2$ par $2x^2 + 6$ sur $\Z$ ?
}

\subsection{Polynômes primitifs}

\nt{
	Alors que dans $\K[x]$, les constantes sont toutes des unités, ce n'est pas le cas dans $\Z[x]$.
	Par conséquent, 2 est un irréductible dans $\Z[x]$ (voir exercice \ref{exe:k-star} et la définition \ref{dfn:irreductible}).
	
	Le polynôme 
		\[ 2x+2 = 2(x+1) \]
	est donc irréductible dans $\Q[x]$ car $2$ est une unité de $\Q$,
	mais n'est pas irréductible dans $\Z[x]$ car $2$ n'est pas une unité de $\Z$.
	
	Il existe donc des polynômes de degré 1 non irréductible dans $\Z[x]$.
	
	Sous condition d'irréductibilité dans $\Q[x]$, nous verrons que l'irréducibilité d'un polynôme de $\Z[x]$ dans $\Z[x]$ est garantie par le fait qu'il ne puisse se factoriser par aucune constante.
	Ces polynômes ont donc des coefficients à PGCD 1, et sont dits {primitifs}.
}

\dfn{polynôme primitif}{
	Un polynôme $f\in\Z[x]$ est un \emphindex{polynôme primitif} dès que
		\begin{align*}
			f(x) = a_0 + a_1 x + \dots a_n x^n, && \pgcd(a_0, \dots, a_n) = 1, && \et && a_n > 0.
		\end{align*}
}{dfn:primitif}

\nt{
	La condition $a_n>0$ permet d'obtenir l'unicité du lemme \ref{lem:primitif} suivant.
	Sans celle-ci, $2(x+1)= -2(-x-1)$ seraient deux décompositions distinctes du même polynôme de $\Z[x]$.
}

\lem{}{
	Chaque polynôme rationnel $f\in\Q[x]$ s'écrit \underline{de façon unique} comme
		\[ f(x) = c \cdot f_0(x), \]
	avec $f_0\in\Z[x]$ primitif et $c\in\Q^\times$.
	
	De plus, $f \in \Z[x] \iff c \in \Z$.
}{lem:primitif}

\nomen{
	La constante $c$ du lemme \ref{lem:primitif} s'appelle le \emphindex{contenu} de $f$.
}

\pf{}{
	Soit $f\in\Q[x]$.
	Posons $d\in\Z^*$ le produit des des coefficients de $f$.
	Alors $g = d \cdot f \in \Z[x]$.
	En factorisant par le PGCD des coefficients de $g$, nous obtenons un $c\in\Q^\times$ tel que $\pm c f$ est primitif, ce qui démontre l'existence.
	
	Supposons désormais que $f$ s'écrive $f = c f_0 = c' f_0'$ avec $c, c'\in\Q^\times$ et $f_0, f_0' \in \Z[x]$ primitifs.
	%Les nombres $c$ et $c'$ sont de même signe car les coefficients dominants de $f_0$ et $f'_0$ sont positifs.
	Nommons $c^{-1}c' = \frac{u}{v}$ avec $u, v\in\Z$. % une représentation sous forme de fraction irréductible.
	Comme $v f_0 = u f_0'$, comparer les PGCD des coefficients donne $v = u$.
	Il suit donc que $f_0 = f_0'$ et $c=c'$.
	
	L'équivalence $f \in \Z[x] \iff c \in \Z$ est claire.
}

\subsection{Classification des irréductibles}

\lem{de Gauss}{\leavevmode
	\begin{enumerate}
		\item
		Le produit de deux polynômes primitifs est primitif.
		\item
		Soit $f_0$ un polynôme primitif et $g \in \Z[x]$.
		Si $f_0|g$ dans $\Q[x]$, alors $f_0|g$ dans $\Z[x]$.
	\end{enumerate}
%	Soit $f\in\Z[x]$ un polynôme à coefficients entiers.
%	D'après le théorème \ref{thm:decomp-Kx}, $f$ se décompose dans $\Q[x]$ de façon unique comme produit de facteurs irréductibles.
%	
%	Si $f$ est unitaire et que les facteurs de la décomposition le sont aussi, alors la décomposition dans $\Q[x]$ est une décomposition dans $\Z[x]$ : tous les facteurs sont unitaires à coefficients entiers.
}{lem:Gauss}

%\nt{
%	Il est possible de décomposer $x^2 - 1$ dans $\Q[x]$ comme
%		\[ x^2 -1 = \left( \frac12 x - \frac12 \right) (2x + 2), \]
%	qui n'est pas une décomposition dans $\Z[x]$,
%	mais les facteurs ne sont pas unitaires ici.
%}

\pf{}{\leavevmode
	\begin{enumerate}
		\item
		Soient $f_0, g_0\in\Z[x]$ deux polynômes primitifs.
		Le produit $f_0g_0$ est à coefficients entiers de coefficient dominant positif comme produit du coefficient dominant de chaque facteur.
		
		Supposons que le PGCD des coefficients du produit ne soit pas $1$.
		Alors un premier $p$ le divise.
		La réduction modulo $p$ donne donc $f_0 g_0 = 0$ dans $\F_p[x]$.
		Or $\F_p[x]$ est intègre d'après l'exercice \ref{exe:Fpx-integre}.
		Il suit donc que $f_0 =0$ ou $g_0 = 0$ dans $\F_p[x]$, ce qui contredit leur caractère primitif.
		\item
		$g = f_0\cdot q$ avec $q\in\Q[x]$.
		Montrons que le contenu de $q$ est égal au contenu de $g$.
		D'après le lemme \ref{lem:primitif}, 
		le contenu d'un polynôme étant dans $\Z$ si et seulement si celui-ci est dans $\Z[x]$, ceci concluera.
		
		Décomposons $g = c_g \cdot g_0$ et $q = c_q \cdot q_0$.
		Alors l'identité
			\[ c_g g_0 = c_q (f_0q_0) \]
		implique que $c_g = c_q$, par unicité du contenu et car $f_0q_0$ est primitif d'après le point 1.
	\end{enumerate}
}

\cor{classification des irréductibles de {$\Z[x]$}}{
	Les polynômes irréductibles de $\Z[x]$ sont exactement, plus ou moins
	\begin{enumerate}
		\item
		les nombres premiers $p$ ; et
		\item
		les polynômes primitifs et irréductibles dans $\Q[x]$.
	\end{enumerate}
}{cor:irr-Zx}

\pf{}{\leavevmode

	\underline{Inclusion $\subseteq$.}
	
	Soit $f\in\Z[x]$ irréductible dans $\Z[x]$.
	\begin{enumerate}
		\item
		Si $f$ est constant, il est nécessairement de la forme $\pm p$ avec $p$ premier, sinon il se scinde en produit de deux facteurs non unités.
		
		\item
		Si $f$ n'est pas constant, la décomposition $f = c \cdot f_0$ du lemme \ref{lem:primitif} n'est une réduction de $f$.
		Dès lors, $c = \pm1$ et $\pm f$ est primitif.
		
		Montrons que $f$ est irréductible dans $\Q[x]$ par contradiction.
		Si $f=ab$ se réduit dans $\Q[x]$, alors les facteurs $a$ et $b$ sont non constants.
		En notant $a = c_a\cdot a_0$ et $b = c_b\cdot b_0$ la décomposition du lemme \ref{lem:primitif}, nous avons
			\[ \pm f = (c_a c_b)\cdot  a_0 b_0. \]
		Le lemme de Gauss \ref{lem:Gauss} implique que $a_0b_0$ est primitif et l'unicité du lemme \ref{lem:primitif} que $c_ac_b = \pm 1$.
		Ainsi $f = \pm a_0 b_0$, décomposition de $f$ dans $\Z[x]$ contredisant l'hypothèse de départ.
	\end{enumerate}
	
	\underline{Inclusion $\supseteq$.}
	
	\begin{enumerate}
		\item
		Soit $f= \pm p$ avec $p$ premier. 
		Alors toute décomposition $f=ab$ implique bien sûr que $a$ ou $b$ est unité.
		Le polynôme $f$ est donc bien irréductible.
		
		\item
		Soit $f\in\Z[x]$ est primitif et irréductible dans $\Q[x]$.
		Si $f$ se réduit dans $\Z[x]$, alors la décomposition $f=ab$ dans $\Z[x]$ contredit l'irréducibilité dans $\Q[x]$ sauf dans le cas où $a$ ou $b$ est constant non unité de $\Z$.
		Dans ce cas, $f$ n'est pas primitif car il se factorise par une constante non unité.
	\end{enumerate}
}

\subsection{L'anneau $\Z[x]$ est factoriel}

\lem{d'Euclide dans {$\Z[x]$}}{
	Le lemme d'Euclide reste vrai dans $\Z[x]$ :
	si $a$ est irréductible et $a|(bc)$, alors $a$ divise $b$ ou $a$ divise $c$.
	
	L'anneau $\Z[x]$ est donc factoriel. Tout polynôme $f \in \Z[x]$ se décompose comme
		\[ f = \pm p_1 \cdot p_2 \cdot p_3 \cdot \cdots \]
	où les $p_i$ sont les irréductibles décrits dans le corollaire \ref{cor:irr-Zx}.
}{lem:Euclide-Zx}

\pf{}{
	Le polynôme $a\in\Z[x]$ est irréductible. 
	La classification \ref{cor:irr-Zx} donne donc les deux alternatives suivantes.
	\begin{enumerate}
		\item
		Si $a = p$ un nombre premier, alors $p|bc$ donne $bc = 0$ dans $\F_p[x]$.
		Par suite, $b =0$ ou $c=0$ dans $\F_p[x]$,
		d'où $p|b$ ou $p|c$.
		\item
		Si $a$ est primitif et irréductible dans $\Q[x]$,
		alors le lemme d'Euclide dit que $a|b$ ou $a|c$ dans $\Q[x]$.
		Les trois polynômes étant à coefficients entiers, et $a$ étant primitif, le lemme de Gauss \ref{lem:Gauss} implique que $a|b$ ou $a|c$ dans $\Z[x]$.
	\end{enumerate}

	S'agissant de l'unicité de la décomposition,
	la preuve du théorème \ref{thm:decomp-Kx} peut être reprise telle quelle.
}

\nt{
	Le lemme d'Euclide donne en fait l'implication « irréductible $\implies$ premier ».
	Une décomposition en facteurs irréductibles est toujours garantie tant qu'il n'existe pas de suite infinie croissante d'idéaux dans l'anneau considéré (c'est-à-dire, l'algorithme d'Euclide aboutit).
	L'unicitié, elle, ne peut qu'uniquement découler du fait que cette décomposition est en facteurs premiers.
}


\section*{Exercices supplémentaires}
\addcontentsline{toc}{section}{Exercices supplémentaires}

\exe{}{
	Vrai ou faux ?
	\begin{enumerate}[itemsep=10pt]
		\item
		Si $p\in\K[x]$ vérifie $p(a) = 0$ pour tout $a\in\K$, alors $p$ est le polynôme nul.
		\item
		La division euclidienne existe dans $\F_{11}[x]$.
		\item
		La division euclidienne existe dans $\Z[x]$.
		\item
		Si $f|g$ dans $\K[x]$, alors $\deg(f) \leq \deg(g)$.
		\item
		Si $f'(x) = 0$ dans $\K[x]$, alors $f$ est une constante.
	\end{enumerate}
}{exe:TF-polynomes}{
	\begin{enumerate}
		\item
		Faux sur les corps finis. Prendre $p(x) = x(x-1)$ dans $\F_2[x]$. 
		\item
		Vrai car $\F_{11}$ est un corps.
		\item
		Faux. La division euclidienne de $x$ par $2$ ne reste pas dans $\Z[x]$.
		\item
		Faux uniquement à cause du polynôme nul $g=0$. 
		Tous les polynômes divisent le polynôme nul. 
		\item
		Faux sur les corps finis. Prendre $f(x) = x^{11}$ dans $\F_{11}[x]$.
	\end{enumerate}
	
}

\exe{difficulty=1}{
	Soient $a, b\in\R[x]$ deux polynômes à coefficients réels.
	Montrer les équivalences suivantes.
	\begin{enumerate}[label=(\roman*)]
		\item $\pgcd(a, b) \neq 1$ dans $\R[x]$
		\item $\pgcd(a, b) \neq 1$ dans $\C[x]$
		\item $a$ et $b$ admettent au moins une racine complexe en commun
	\end{enumerate}
}{exe:gcd-roots-R-C}{

	\underline{Implication (i) $\implies$ (ii).}
	
	Si $d\in\R[x]$ non 1 divise $a$ et $b$ dans $\R[x]$, alors $d$ divise également $a$ et $b$ dans $\C[x]$.
	
	\underline{Implication (ii) $\implies$ (iii).}
	
	Soit $d=\pgcd(a,b) \in \C[x]$. Comme $d$ est non constant par hypothèse, $d$ admet une racine complexe $z\in\C$ d'après le théorème fondamental de l'algèbre.
	Ainsi $d(z) = 0$ et donc $a(z) = 0$ et $b(z) = 0$ car $d$ divise $a$ et $b$.
	
	\underline{Implication (iii) $\implies$ (i).}
	
	Notons $z\in\C$ une racine complexe en commun à $a$ et $b$.
	Par contradiction, si $\pgcd(a, b) = 1$ dans $\R[x]$, le théorème de Bézout impliquerait l'existence de polynômes $u, v \in\R[x]$ tels que $au + bv = 1$.
	En plongeant cette égalité dans $\C[x]$, évaluer en $x=z$ aboutit à la contradiction $0=1$. \Lightning
}


\exe{difficulty=1}{
	Soit $z\in\C$ un \emphindex{nombre algébrique}, c'est-à-dire une racine d'un polynôme sur $\Q[x]$.
	\begin{enumerate}
		\item Montrer que l'ensemble 
			\[ N = \bigset{ f \in \Q[x] \tq f(z) = 0 } \]
		est un idéal de $\Q[x]$.
		\item
		Soit $m_z \in \Q[x]$ le \emphindex{polynôme minimal annulateur} de $z$ : c'est le polynôme de $\Q[x]$ annulant $z$ et de degré minimal.
		Montrer que $m_z$ est bien unique et qu'il est irréductible dans $\Q[x]$.
		\item
		Montrer que $N = (m_z)$. C'est-à-dire que tout polynôme annulateur de $z$ est divisible par $m_z$ dans $\Q[x]$.
	\end{enumerate}
}{exe:poly-min}{
	todo
}

\exe{difficulty=1}{
	Considérons $\Q[\sqrt2] = \bigset{ r + s\sqrt2 \tq r, s \in \Q}$ ainsi que l'application \mbox{« évaluation en $\sqrt2$ »}
		\begin{center}
		$\ev_{\sqrt2}$ :
		\begin{tabular}{ccc}
			$\Q[x]$ & $\longrightarrow$ & $\Q\left[\sqrt2\right]$ \\
			$f$ & $\longmapsto$ & $f\bigpar{\sqrt2}$
		\end{tabular}
		\end{center}
	Montrer que $\ker(\ev_{\sqrt2}) = ( x^2 - 2 )$ et que l'idéal est un \emphindex{idéal maximal} : il ne contient pas de sous-idéal propre qui n'est pas $\Q[x]$ tout entier.
}{exe:Qsqrt2}{
	todo
}

\exe{}{
	Soit $p\in\K[x]$ un polynôme premier, c'est-à-dire non nul et vérifiant, pour tout $a, b\in\K[x]$, $p|(ab) \implies p|a$ ou $p|b$.
	Montrer que $p$ est irréductible.
}{exe:prime-implies-irr}{

}

\exe{}{
	Soit $p\in\K[x]$ un polynôme.
	Montrer que $p$ est irréductible si et seulement si $p(x+c)$ est irréductible, où $c\in\K$ est un nombre quelconque.
}{exe:irr-shift}{
	todo
}

\exe{difficulty=2}{
	Montrer le \emphindex{critère d'Eisenstein} :
	soient $f(x) = a_0 + a_1x + \cdots + a_nx^n \in\Z[x]$ et $p$ un nombre premier vérifiant les conditions
	\begin{enumerate}
		\item $p$ ne divise pas $a_n$ ; et
		\item $p$ divise $a_0, a_1, \hdots, a_{n-1}$ ; et
		\item $p^2$ ne divise pas $a_0$.
	\end{enumerate}
	Alors $f(x)$ est irréductible dans $\Q[x]$.
}{exe:eisenstein}{
	Comme $f$ est à coefficients entiers, le lemme de Gauss \ref{lem:Gauss} implique que l'irréducibilité dans $\Q[x]$ est équivalente à l'irréducibilité dans $\Z[x]$.
	Supposons donc, par contradiction, que $f$ soit réductible sur $\Z[x]$.
	
	Alors $f = ab$ avec $a, b\Z[x]$ non constants de coefficients dominants ne divisant pas $p$.
	En prenant la relation modulo $p$, le polynôme $f$ devient réductible sur $\F_p[x]$.
	Or, par hypothèses sur les coefficients de $f$, nous avons que $f = a_n x^n \mod p$ avec $a_n$ une unité de $\F_p$.
	Comme $x$ est irréductible dans $\F_p[x]$, la seule manière de réduire $f$ dans $\F_p[x]$ est de le décomposer comme $a_n \cdot x^k \cdot x^{n-k}$.
	Ainsi, les deux facteurs $a$ et $b$ sont congrus à un multiple d'une puissance de $x$ modulo $p$, non constante car $a$ et $b$ sont unitaires non constants.
	Il suit donc que $p$ divise le coefficient constant $a(0)$ et $b(0)$, et donc que $p^2$ divise le coefficient constant $a_0 = f(0) = a(0)b(0)$.
	Ceci contredit la troisième condition imposée à $f$. \Lightning
}

\exe{}{
	Montrer que le polynôme $f(x) = x^5 - 10x^3 + 50x - 20$ est irréductible dans $\Q[x]$.
}{exe:Eisenstein-example}{
	Le critère d'Eisenstein de l'exercice \ref{exe:eisenstein} avec $p=5$ permet de conclure.
}

\exe{difficulty=1}{
	Considérons la fonction rationnelle $f(x) = \frac{x^p - 1}{x-1}$ avec $p$ un premier.
	\begin{enumerate}
		\item Montrer que $f$ appartient à $\Z[x]$.
		\item Montrer que $f$ est irréductible en montrant que $g(x) = f(x+1)$ vérifie le critère d'Eisenstein de l'exercice \ref{exe:eisenstein}.
	\end{enumerate}
}{exe:xppm1-irr}{
	todo
}

\exe{difficulty=1}{
	Soit $n\in\N$ et $p\neq2$ un nombre premier.
	Le but de cet exercice est de démontrer que le polynôme $f(x) = x^n + x + p$ est irréductible dans $\Q[x]$.
	\begin{enumerate}
		\item
		Justifier qu'il suffit de démontrer que $f$ est irréductible dans $\Z[x]$.
		\item
		En voyant $f$ comme polynôme de $\C[x]$, montrer que toutes ses racines $z\in\C$ sont de module strictement supérieur à 1.
		\item
		En déduire que si $f=ab$ dans $\Z[x]$, alors $|a(0)|, |b(0)| > 1$.
		\item
		Conclure.
	\end{enumerate}
}{exe:xn-n-p-irr}{
	\begin{enumerate}
		\item
		Si $f$ est réductible dans $\Q[x]$, alors, comme $f$ est unitaire de $\Z[x]$, le lemme de Gauss (polynômes) \ref{lem:Gauss} garantit que $f$ est réductible dans $\Z[x]$.
		\item
		Soit $z\in\C$ tel que $f(z) = 0$.
		Alors $z^n + z + p = 0$, et donc $2 < |p| \leq |z| + |z|^n$.
		Il suit donc naturellement que $|z|>1$ car, dans le cas contraire, nous aurions la contradiction $2 < |z|+|z|^n \leq 2$.
		\item
		Le coefficient constant de $a$ vu comme polynôme de $\C[x]$, $a(0)$, est égal au produit de ses racines.
		Or, les racines de $a$ sont exactement des racines de $f$, dont le module est strictement supérieur à 1.
		Il suit donc que $|a(0)| > 1$.
		Le cas pour $b$ est identitque.
		\item
		Supposons que $f$ se réduise comme $f=ab$ dans $\Z[x]$.
		Alors $f(0) = p = a(0)\cdot b(0)$, égalité dans $\Z$.
		Comme $p$ est premier, soit $a(0)$, soit $b(0)$ est une unité de $\Z$, ce qui contredit que $|a(0)|, |b(0)| > 1$.
	\end{enumerate}
}


\shipoutExercise

\section*{Solutions}
\addcontentsline{toc}{section}{Solutions}

\shipoutAnswer



